\chapter{Intelligence Inflation---The Arrival of the Affluent Society}


\begin{myquote}
The future has arrived, and it is distributed more evenly than ever before. Power belongs to everyone.

\hfill ---Andrej Karpathy
\end{myquote}

{\kaiti
A friend told me he used AI to write a love letter to his girlfriend.

``It was really good---ten times better than anything I could write myself. She cried when she read it.''

I asked: ``Does she know AI wrote it?\,''

``No.''

I said: \textbf{``Then you're doomed.''}
}
\vspace{1em}

He didn't understand: the result was better, the process more efficient, and she was more moved---what's the problem?

\textbf{The problem is enormous.}

There are two fundamentally different situations when AI does things for us: one is augmentation, the other is forgery. When AI simultaneously analyzes ten thousand financial reports and discovers patterns invisible to the human eye---that is augmentation.

\vspace{1em}

But a love letter is different.

What is the value of a love letter? Not the words themselves, but that I was willing to spend two hours for you, racking my brain, writing down these clumsy sentences. It is those two hours, that ``racking my brain,'' that clumsiness.
AI writes more fluently, more beautifully, more movingly---but that ``willingness'' has been stolen.
His girlfriend cried, not because the words were beautiful, but because she believed he had poured his heart into them for her.
What moved her was the devotion, not the quality. An AI-written love letter means that devotion never existed at all.

\vspace{1em}

Time, clumsiness, caring.
These things that cannot be replicated may be the most precious things in the future.

%========================================================
\section{Distributed More Evenly Than Ever: What AI Brings Is Not Shared Prosperity}
%========================================================

Among the many discussions about AI's impact on society, Andrej Karpathy flipped William Gibson's famous line---``The future is already here---just not evenly distributed.''

He noticed something that doesn't quite follow the usual path of technology diffusion: AI has not spread along the traditional routes. Many previous revolutionary technologies---nuclear energy, spaceflight, the internet---were typically controlled first by governments or large institutions, then gradually trickled down to the public.

\vspace{1em}

This time, the order is reversed. His observation is not complicated, yet it is crucial: in many fields, ordinary people have adopted AI faster than institutions. Employees often start using it first; organizations only slowly acknowledge, regulate, and restructure their processes around it. AI was not first locked away in a handful of institutions and then released layer by layer. It quickly landed on the devices that vast numbers of ordinary people already owned.

What does this mean?

Compared to many previous technological revolutions, this time ordinary people need not wait long for the trickle-down effect. A twenty-year-old college student and a fifty-year-old CEO face the same tool; an entrepreneur in Lagos and an engineer in Silicon Valley may be calling the same model. This does not mean inequality has disappeared---far from it---but in quite a few use cases, the gap at the starting line has been noticeably compressed.

\vspace{1em}

Karpathy admitted he has ``never felt so behind.''

He described AI tools this way: ``A powerful alien artifact has been distributed to everyone, but it came with no manual.'' Then he added the second half: ``If I can string these tools together correctly, I can become ten times more powerful.''
The disruption is real, and so is the opportunity---open to all.

\vspace{1em}

The lever has been distributed to everyone, but the ability to use the lever will not be evenly distributed along with it.

When the barrier to tools drops, what truly separates people is no longer just resources, but judgment, taste, and the willingness to take risks. When the gap at the starting line narrows, divergence may actually accelerate. Because leverage amplifies not only capability, but the differences between capabilities.

This is the affluent society: not evenly shared prosperity, but swifter divergence once the gates open to all. Nearly everyone has their hand on the same door handle, but the moment the door swings open, divergence begins again.

%========================================================
\section{The Contours of Affluence: New Scarcities Will Sprout in Unexpected Places}
%========================================================

Over the past two years, some of Silicon Valley's most influential technology optimists have been painting the same picture over and over: Sam Altman emphasizes that the cost of intelligence will keep plummeting; Dario Amodei pushes this vision to its extreme, depicting a ``nation of geniuses in data centers''---countless intelligent agents working continuously at an efficiency far beyond human capacity; Jensen Huang takes it a step further, reminding us that this need not be a zero-sum substitution, because new intelligent systems will also create new possibilities.

\vspace{1em}

If you listen only to these voices, the future might be summed up in a single word: affluence.

\vspace{1em}

This optimism is not entirely without basis. Erik Brynjolfsson of Stanford University and others have provided more concrete evidence: in a customer support setting at a large software company, after integrating a generative AI conversational assistant, agents resolved an average of roughly 14\% more issues per hour, with new hires and lower-skilled employees seeing productivity gains of approximately 34\%\footnote{Brynjolfsson, Li, Raymond, \textit{Generative AI at Work}, NBER Working Paper 31161, 2023.}. This suggests that AI is not merely an amplifier for elites---in certain scenarios, it can also serve as an equalizer for ordinary people.

\vspace{1em}

But affluence does not mean universal access.

Technology may continue to drive down the unit cost of intelligence, but for ordinary people, the real question is not ``Is there cheaper intelligence somewhere in the world?'' but rather ``Can I afford to keep using it?''

\vspace{1em}

Free-tier, queued, and ad-supported AI may become as ubiquitous as tap water, but truly stable, reliable, always-on, high-quality intelligence that can continuously work on our behalf may---like today's private physicians, top-tier lawyers, and first-class cabins---become yet another tiered service.

\vspace{1em}

The question is not whether intelligence will become cheap, but rather: once human income has been repriced, who can still afford ``good enough intelligence''?
Intelligence may become abundant, yet scarcity never disappears---it only migrates. Every time an old scarcity is eliminated, a new one sprouts in an unexpected place. And we will reorganize our social hierarchies around these new scarcities.

\vspace{1em}

So along what path will affluence unfold?
It extends along at least three fault lines---the leverage of the individual will be pushed to its limit, the form of organizations will undergo a phase transition, and the meaning of money will be rewritten.

\newpage
%========================================================
\subsection{The Leverage of the Individual: Everyone Will Have an Always-On Agent}
%========================================================

Throughout human history, leverage has come in three forms\footnote{Naval Ravikant's leverage framework.}:

\textbf{Labor}---hiring people, the oldest and most constrained;

\textbf{Capital}---making money with money, powerful but requiring accumulation;

\textbf{Code and media}---create once, replicate infinitely, marginal cost zero.

The third is ``leverage without accumulation''---write a piece of software, post a tweet, no one's permission required. AI has pushed this third form of leverage to limits previously unimaginable.

\vspace{1em}

The venture capital industry now treats ``a ten-person, hundred-billion-dollar company'' as a serious hypothesis. It is no longer a science-fiction trope; it is redefining how people imagine the scale of organizations.

\vspace{1em}

Peter Steinberger may be the best living footnote to this hypothesis.

This Austrian spent more than a decade building PSPDFKit, a company that processes PDFs. After the tools he created came to be used by vast numbers of developers worldwide, he completely burned out---``I stared at the screen and couldn't write a thing. My mind was blank.'' He bought a one-way ticket to Madrid and disappeared.

\vspace{1em}

In the spring of 2025, he stumbled upon the fact that the world had changed---AI could now handle those repetitive coding tasks, so he started building projects on his own. OpenClaw was his forty-fourth AI project; the first forty-three had all failed.
His single-day code commit count on GitHub exceeded a thousand. This was not the upper limit of human typing speed---it was the result of him driving multiple AI Agents in parallel.

He didn't need to write every line of code himself---he was responsible for defining tests, reviewing logic, and making architectural decisions, while the AI Agents handled concurrent execution. He said: ``I don't read code anymore. I just watch the code flow past.''

This is the new landscape once leverage is pushed to its limit: a single person no longer just does one thing faster, but gains the ability to orchestrate parallelism that once required an entire organization.

\vspace{1em}

Steinberger's leverage operates in two directions simultaneously.

The first direction is vertical---amplification of a single capability. AI enables one person to produce the volume of code that previously required a whole team. This is the kind of leverage most people think of when they talk about AI: faster, more, stronger.

\vspace{1em}

\newpage
The second direction is horizontal---expansion of parallel threads. Steinberger is not running faster on a single track; he is advancing along multiple tracks at once. He maintains several projects simultaneously, switches between multiple contexts, and makes ``good enough'' judgments amid imperfect information. Like an orchestra conductor---he need not play any instrument better than the musicians themselves, but he must hear all the instruments at once and make judgments at the level of the whole.

\vspace{1em}

Vertical leverage is something AI can directly provide---it is the tool that makes us stronger at any single point. But horizontal leverage---how many threads you can orchestrate simultaneously, how many parallel contexts you can sustain judgment across without collapsing---depends on the person, not the tool. This illustrates the essence of the ``super individual'' more vividly than any sermon: \textbf{not a single stroke of genius, but high bandwidth plus high fault tolerance.} Being able to absorb forty-three failures and still keep swinging is itself a form of bandwidth.

\vspace{1em}

The product OpenClaw itself pushes the meaning of leverage one layer further.

When mobile internet emerged, there was a precise metaphor: the phone is the outlet, and the person is the plug.

The outlet is always on---power, network, signal, available at all times.

But the person, as the plug, can be pulled out---turn off the phone, go to sleep, go stare into space.

The device is always on; the person can go offline.

\vspace{1em}

Now everyone may have an Agent that is always on. It doesn't need to sleep the way humans do; it can negotiate car prices for us after we fall asleep, handle emails, monitor data, and in an ever-growing number of routine matters, act on our behalf.

\vspace{1em}

In the mobile internet era, the bottleneck of leverage was how much time we were willing to spend staring at a screen. In the AI era, the bottleneck has shifted---it is no longer how much time you can invest, but how many threads you can orchestrate at once. This is the ``bandwidth'' of leverage.

\vspace{1em}

When an individual can do what once required an entire company, does the company as an organizational form still need to exist?

%========================================================
\subsection{The Phase Transition of Organizations: The Company May No Longer Be the Most Efficient Form}
%========================================================

In 1937, the young economist Ronald Coase asked a question that seemed almost naive: if markets are so efficient, why do firms exist at all?

His answer became a cornerstone of organizational theory: firms exist because market transactions have costs. Finding collaborators takes time, negotiating contracts takes effort, monitoring execution takes resources. When these ``transaction costs'' exceed the management costs within a firm, it becomes more economical to hire people into the same organization and coordinate through hierarchy than to outsource every task on the open market.

\vspace{1em}

The boundary of the firm is drawn at the line where transaction costs and management costs are equal.

AI is pushing that line to an extreme. When Agents can complete information searches, contract drafting, credit verification, and task allocation in milliseconds, the friction of market transactions drops precipitously. If market transaction costs approach zero, part of the economic rationale for the firm's existence is undermined.

\vspace{1em}

% But the software world may have exposed another layer of truth first: as execution itself becomes ever cheaper, the value humans provide no longer resides primarily in ``making things with their own hands,'' but migrates to another layer---one far harder to outsource.

% Code will get cheaper and cheaper; ground truth will get more and more expensive. Systems can endlessly generate, test, patch, and deploy, but they will not answer the more costly questions on their own: Is this thing actually worth doing? What counts as doing it right? What sacrifices are acceptable? If something goes wrong, who bears the consequences? Once execution is automated, the most expensive thing is no longer labor, but judgment and accountability.

Some large institutions have already coined a name for this new form: the frontier company---powered by on-demand intelligence, operated by hybrid teams of humans and AI Agents, organized around output targets rather than functional departments. The new role of employees resembles that of an ``Agent boss''---humans define objectives, assign tasks, and bear responsibility; Agents execute. The flip side is equally visible: organizations will become more agile, and more ruthless. Efficiency gains and layoff pressures often arrive together.

\vspace{1em}

So does the company still need to exist? Perhaps, but the reasons are undergoing a phase transition.
Part of its productive function will be weakened, while its role as a container of liability, a node of authorization, and a bearer of consequences may become even more important.
What we may witness is not the death of the firm, but its reorganization---in some domains, alliances, studios, micro-companies, and protocol-based entities will emerge, while heavily regulated, liability-laden domains will preserve the shell of the large institution.
The ``organization'' of the future may split into two forms: one lightweight and protocol-driven, a production network; the other heavy and identity-driven, a container of belonging.
The two may look nothing alike, or they may be two sides of the same coin.

% Transaction costs are only one reason firms exist. The firm is also a governance device, a capital container, and a liability vehicle. Even if agents make assembling ``micro-organizations'' cheap, the expensive part may remain accountability: who has the authority to authorize action? Who bears legal risk? When an autonomous workflow causes harm, who can be held responsible? In a world where algorithms can draft contracts but cannot sign on the dotted line, the concept of the ``legal person'' will not vanish---it will be refilled.

% Although the technology is ready, organizational change often lags behind. Many large companies recognize AI's importance, but their management transformation has stalled. Back-end architectures built around ``human workforce management'' cannot easily support the productivity leaps driven by ``compute power.'' The organizational inertia of legacy systems is more stubborn than most people imagine.

% At the same time, the psychological function of the company will not be erased by the erosion of economic logic. ``I work at such-and-such company''---this sentence provides not only career information but social coordinates, group identity, and an anchor for one's life narrative. When the economic logic of the firm weakens, its psychological function may actually become more prominent---we join an organization not because market transaction costs are too high, but because the cost of loneliness is too high.

% So ``phase transition'' may be a more accurate metaphor than ``replacement.'' What we may see is not the death of the firm but its reorganization---alliances, studios, micro-companies, and protocol-based entities in some domains, with heavily regulated and liability-laden domains preserving the shell of the large institution. The ``organization'' of the future may split into two forms: one lightweight, protocol-driven production network, the other a heavy, identity-driven container of belonging. They may look completely different. Or they may be two sides of the same thing.

% Klarna's story may be the most instructive case in this phase transition. The Swedish fintech company froze hiring, cut its workforce from five thousand five hundred to three thousand, and claimed its AI assistant could do the work of seven hundred customer service agents. The CEO used his own AI clone to deliver quarterly earnings. Then customer satisfaction plunged. The CEO's tone changed completely: ``Cost unfortunately became an overly dominant evaluation factor. The result was that quality declined.'' He began rehiring human agents.

% Turning ice into water takes time. Trying to skip the process---burning ice directly with fire---yields not water, but cracks.
\vspace{1em}
Organizations are being reshaped; individuals are expanding.

But what about the medium of exchange that runs through all of this---money?

If the cost structure of production has changed, what happens to the role of money?


\newpage
%========================================================
\subsection{The Role of Money: IOUs Are Increasingly Made Out to the Owners of AI}
%========================================================
% But things may not be so simple. Perhaps the more worthwhile question is not whether the cost of compute will drop, but: even if it does, who will pay for our share of that compute?

% No matter how cheap compute becomes, it is still a cost. Someone has to bear that cost.

% The real issue here is that money still anchors to labor, but the labor humans can provide, relative to the total supply, is low in quality and unstable in availability.

% So it's not that money is useless, but that ordinary humans cannot provide the corresponding anchor.

% And while the cost of compute may decrease, the selling price does not necessarily decrease---especially when you have nothing of value to exchange.


The most elemental understanding of money is this: an IOU for labor.

Person A does something for Person B; Person B gives Person A a token, redeemable for someone else's labor.

What anchors it, at the bottom, is our labor and our time.

Elon Musk pushed this logic to its extreme---money is nothing more than a database for allocating labor.

If you had fully automated robots with unlimited energy, money would lose its reason to exist.

\vspace{1em}

So will money become obsolete? \textbf{Probably not.}

Even if AI makes the supply of labor overwhelmingly abundant and its cost vanishingly low, money can still measure the exchange of value---only the composition of labor shifts from purely human supply to humans plus AI.
Money, as a bookkeeping system, still works.

\vspace{1em}

The second question has a bigger impact on ordinary people: who is the IOU made out to?

The volume of translation work that a human translator could handle in a day can now be drafted by a model in a fraction of the time, at a fraction of the cost. The translation market hasn't disappeared---the total volume has actually grown.

But the supply side has changed: AI accounts for the overwhelming majority; humans are left with a sliver.

\vspace{1em}

Money still circulates, but the share of human translators' labor in total supply has gone from a meaningful figure to a negligible rounding error.
Money hasn't stopped working; what has stopped working is the portion of our labor we used to exchange for it.
IOUs are still being written, but increasingly they are made out to the owners of AI.

\vspace{1em}

So the fundamental question is not whether compute will become cheaper, but rather: once human labor has been repriced, who can still produce something worth exchanging for it?
Declining costs do not mean everyone can afford it---``cheap'' is a word that applies only to those who still have something to trade. For those squeezed out of the exchange, the same compute may still be a barrier: when our labor amounts to nothing more than a rounding error in total supply, on what basis do we claim the right to use that compute? On the basis of institutions? Of a social contract?

\vspace{1em}

These are all possible answers, but none of them happens automatically.

Money has not become useless. The IOUs are still there; it is just that the names on them are changing.

\newpage
%========================================================
\section{The Usefulness of the Useless: Becoming What Lies Outside the Distribution}
%========================================================

We have grown accustomed to an equation: useful = capable of generating productivity.

This equation felt so natural in industrial society that we forgot it was man-made. A person's value was defined by output---how many garments produced, how many lines of code written, how many spreadsheets processed. Because making a living required compliance, individuality was the enemy of productive efficiency. On the assembly line---whether physical or cognitive---standardization was profit.

\vspace{1em}

Then AI arrived, better at standardized tasks than any standardized human.
If AI becomes the primary source of productivity, then \textbf{the definition of ``useful'' will need to be rewritten}.

Perhaps the answer lies in the ``useless.'' ``The usefulness of the useless'' comes from Zhuangzi, and today it takes on new meaning. When the marginal cost of writing a report, building a website, or translating a book approaches zero, ``useful for material production'' is no longer the sole yardstick for measuring a person's worth.

\vspace{1em}

In statistics, the center of the normal distribution is the most crowded region---most data points cluster around the mean. In the past, industrial production demanded standardized parts. Today, if AI can handle all the work ``within the distribution''---those standardized, modelable tasks that have appeared in training data---then the scarce value we can offer one another likely comes only from patterns AI has never seen, from what lies statistically \textbf{outside the distribution}. We can even foresee that as long as social structures exist, differentiation between people will exist. This has nothing to do with whether AI exists: if we stand ``one leaf'' taller than those around us, if we are ``different,'' then we have unique value.

\vspace{1em}

Outside the distribution means scarce, and scarce means valuable.
The core is still exchange---we must still offer each other something scarce.
Only the definition of ``scarce'' has changed: no longer ``able to produce material goods that others cannot,'' but ``able to provide experiences, perspectives, or connections that others cannot.''

If we are no longer forced to crowd into the center of the distribution as standardized parts, but instead scatter in our own unique directions, the world will grow wider boundaries again. The next generation of AI will of course learn from this distribution that humans have stretched anew, but the flywheel's starting point is still us: someone must first take the risk, make the mistake, pay the irreversible price, before machines can know where the new boundaries lie.

So what is truly worth anticipating about the affluent society is not merely that machines have learned more, but that we finally have the chance to stop living so much like machines---and that we can keep pushing ``outside the distribution'' ever further.

I paused for a long time when I reached this point.

This book is written for my daughter, with an implicit reader in mind: someone young, educated, and who still has time to adapt. ``Become someone outside the distribution,'' ``train your own objective function,'' ``stay curious.''
I still believe in this advice, but it rests on a premise I have never stated plainly---\textbf{that you still have margin.}

\vspace{1em}

Some people have no margin left.

For a fifty-year-old truck driver, autonomous driving is encroaching on his livelihood. Telling him to ``become a super individual'' is cruel---what he needs is not taste and an objective function, but a job tomorrow. A thirty-five-year-old translator discovers that clients, after using AI to produce a first draft, are only willing to pay proofreading rates. Her decade of accumulated expertise has not become ``private context''---it has simply been zeroed out. A forty-year-old graphic designer whose aesthetic sensibility was once a source of pride now hears Midjourney users say ``close enough is fine.'' They don't even know that my daily work is building the very intelligence that replaces them.

\vspace{1em}

They are not lacking in effort, not lacking in curiosity, not lacking an objective function.

They simply found themselves standing on the fault line of transformation, and fault lines do not choose their victims.

\vspace{1em}

In every technological revolution in history, there has been a group of people like this.

When the steam engine arrived, textile workers smashed machines; when the automobile arrived, coachmen took to the streets in protest. Posterity always says ``they should have embraced change.'' But those who say this are usually standing on the flat ground that exists only after the change is complete. People standing at the edge of a cliff cannot hear that kind of advice.

\vspace{1em}

I do not have the ability to solve this problem, nor can this book solve it. Institutional design, social safety nets, wealth redistribution---these are matters of policy, requiring people with greater power to advance them.

\vspace{1em}

There are only two things I can do.

First, not pretend these people do not exist.

Second, if any reader of this book someday walks into the room where policy is made, please remember the people who could not be written into this book. \textbf{The beneficiaries of technological change have an obligation to attend to those who bear its costs}---not out of kindness, but because a narrative with only winners will eventually be shattered by reality.

\newpage
%========================================================
\section{Gears at Different Speeds: One Era, Multiple Tempos}
%========================================================

In Silicon Valley's AI labs, the clock is measured in weeks.
Every week, models grow stronger, benchmarks climb, revenue curves rise at exponential slopes.
The people in the labs talk about ``a nation of geniuses'' as though discussing next year's product roadmap.
They are at 100x speed.

\vspace{1em}

AI practitioners and engineers, product managers, and designers in the tech industry---those who use AI tools daily and see their efficiency multiply---they are at 10x speed. The feeling of being ``ten times more powerful'' that Karpathy described captures exactly the experience of this tier.

\vspace{1em}

Then there is everyone else. In the corridors of pharmaceutical companies, candidate molecules designed by AI in days still await the slow grind of clinical approval. In sub-Saharan Africa, the rollout of certain well-established vaccines remains constrained by infrastructure and governance capacity. They are at 1x speed.

\vspace{1em}

100x, 10x, 1x---these numbers are rhetorical; don't take them literally.

They are like artifacts from different eras that happen to share the same calendar.

Is the resistance to diffusion large or small? There are two views, and both may hold true simultaneously.

\vspace{1em}

The first holds that diffusion will be extremely fast. The argument is straightforward: AI's medium is software, and the marginal cost of distributing software is zero.
Once a model is trained, it can be rapidly distributed to vast numbers of people around the world who have a network connection.
Moreover, AI is a self-accelerating technology: it makes developers more efficient, more efficient developers build better AI tools, better AI tools turn more people into developers, and so the cycle repeats.

\vspace{1em}

The second holds that diffusion will be slow in many domains. The argument is equally straightforward: as long as any single link in the chain must be completed offline, the speed advantage of the digital world will be dragged down by the friction of the physical one.

\vspace{1em}

Tyler Cowen's distinction may come closest to reality: domains that can be completed entirely online with relatively little regulation will transform rapidly; the transformation of government, healthcare, and education will likely take thirty years.

\newpage
%========================================================
\section{The Shape of the Transition: Has an Earthquake Struck Beneath the Calm Surface?}
%========================================================

Regardless of speed, what must come will come. What does this transitional period look like?

Economics offers at least two reference frames:

\textbf{The Solow paradox:} We can see AI everywhere except in the productivity statistics---a transformative technology visible in every domain, yet invisible in the very metrics designed to measure its impact.

\textbf{Engels' pause:} ``Growth'' and ``broadly shared improvement in welfare'' can decouple for extended periods.

Different frameworks point to the same conclusion: patience. The gains will come, but they require waiting.

\vspace{1em}

Historical technology transitions have repeatedly confirmed this rhythm. The spinning jenny, the story of horse-drawn carriages giving way to automobiles, and so on---these were all technologies that automated physical labor, machines replacing muscle.
\vspace{1em}

Our own history offers another mirror---driven by institutional reform, yet equally instructive.

From the late last century into the early years of this one, China underwent a wave of drastic restructuring euphemistically termed ``transitional unemployment.'' Ordinary people called it what it was: layoffs. Its reach was vast; countless families were swept up in it.

\vspace{1em}

This round of reform is often viewed as one of the key preconditions for China's entry into the WTO and the subsequent era of rapid growth. But for the individuals involved, macroeconomic gains did not automatically flow back to those who paid the price, much less arrive in the same year.
The costs were borne by one generation; the benefits only began to materialize in the next.

\vspace{1em}

These two kinds of precedent---technology-driven and institution-driven---share the same premise: there was always some territory that machines could not reach. The spinning jenny replaced the hand, electricity replaced the horse, institutional restructuring eliminated inefficient positions, but cognitive ability remained humanity's last comparative advantage.

\vspace{1em}

AI eliminates that fallback---and this is unprecedented.

No historical precedent can fully map onto this transition.

\vspace{1em}

Earthquakes have already struck in individual industries, but the shockwaves have not yet propagated to macroeconomic indicators.

It is as though the earthquake occurred on the ocean floor---fishing boats near the epicenter have already capsized, but the tide gauges on shore still show steady readings.

The hallmark of a tsunami is this: by the time people on shore notice the water level changing, it is already too late to move to higher ground.

\newpage
%========================================================
\section{The Resistance of Diffusion: How Much Time Do We Have Left to Adapt}
%========================================================
In discussions of technology, the word ``diffusion'' often carries a latent pejorative connotation.

It implies that the thing has already been built---it just hasn't spread yet.

As though resistance itself were a defect, an inefficiency that needs to be eliminated.

\vspace{1em}

Perhaps we should look at friction differently:

A tree does not grow faster just because we give it more information;

% A child does not skip adolescence just because they have a better AI teacher.

% A clinical trial cannot skip the control group just because AI predicts the drug will work.

A society does not automatically become ready to accept a technology simply because the technology is mature.

\vspace{1em}

The physical world has its own rhythm.

Molecules need time to fold, organizations need time to restructure, trust needs time to accumulate.

A legislator needs time to understand what, exactly, the thing they are legislating actually is.

Not all of these are defects; some of them are protective mechanisms.

\vspace{1em}

If diffusion were infinitely fast, if every technological breakthrough permeated every corner of society the instant it was born, we would lose all buffer zones---no time to assess consequences, no time to correct errors, no time for society's immune system to respond.

\vspace{1em}

Part of the reason human civilization was able to digest technological transformations like the Industrial Revolution and nuclear technology is precisely that diffusion was not infinitely fast. We had decades to develop labor unions, to refine nuclear non-proliferation treaties and privacy legislation.
These institutions are imperfect, but their very existence depends on a premise: that the social consequences of technological progress unfold at a speed to which humans can react.

\vspace{1em}

Now, that premise is being challenged.
Not because diffusion has become infinitely fast---even the CEOs of AI companies acknowledge it won't---but because it has become faster than any technology before.

\vspace{1em}

The narrow zone between ``faster than ever before'' and ``not infinitely fast'' may be exactly the window within which our generation must complete all the critical institution-building.

\newpage
%========================================================
\section{Coda: In the Affluent Society, the Scarcest Thing Is Showing Up in Person}
%========================================================

When we think about AI, we often carry a kind of latent anxiety---

Will it help us, or will it replace us?

Perhaps the question itself deserves reexamination.

\vspace{1em}

Ordinary people are already answering these questions in their own way.

Some have discovered that their parents use AI assistants far more than expected.

Not to boost productivity, but to get through the day---

Setting medication reminders, having it keep them company in conversation, treating it as a life assistant available on demand.

``More reliable than your father---it doesn't forget.''

They don't understand large language models; they don't understand prompting.

But they understand one thing better than anyone: technology exists to serve people.

Perhaps this is what ``augmentation'' looks like in its most elemental form---

Not a 10x lever, not a hundred-billion-dollar company, but an elderly person who took their medication on time because AI reminded them.

\vspace{1em}

Karpathy said: ``I've never felt so behind.''

Then he added: ``I can become 10 times more powerful.''

The distance between those two sentences is the distance we are crossing right now.

% And we will have two reactions.

% One is fear---freezing in the face of the unknown, a response that once saved our lives on the savanna.

% The other is curiosity---reaching out to touch the unknown, the response that once led us off the savanna.

% Both reactions are reasonable, but curiosity is more useful than fear.

% Not because there is no danger---the danger is real.

% But because fear makes us stop, and curiosity makes us walk toward it.

\vspace{1em}


% Anyone who claims to know the final outcome is either selling something or hasn't been paying close enough attention. The honest position is: maintain fundamental uncertainty about the endgame, while being highly confident that the transition itself will be difficult, uneven, and will reshape human identity in ways we are only beginning to articulate.


But perhaps the endgame matters less than we think.

What we are living through is not an endpoint that needs to be predicted,

but a passage that needs to be traversed.

Along this passage, some things will become cheap;

others will become immeasurably precious---time, clumsiness, caring.

\vspace{1em}

The friend who used AI to write a love letter---his mistake was not in using a tool. His mistake was that he himself was not present.
In all this affluence, the scarcest thing may be a person's \textbf{willingness to do something for another person---clumsily, inefficiently, and in person.}
